% GENERATED FILE. Edit canonical lessons, then run npm run generate:books.
% canonical-source-hash: fnv1a64:278e3a4d3ac56eb7
% canonical-lessons: GE-C447-friseur, GE-C447-foehnen, GE-C447-optiker, GE-C447-uhrmacher, GE-C447-juwelier

\chapter{Friseur und Optiker}
\label{ch:ge-friseur-und-optiker}
\hlchaptermodality{447}

\begin{chapteropening}
\textbf{By the end of this chapter:} \emph{I can use personal services: the hairdresser, blow-drying, the optician, the watchmaker and the jeweller.}

Complete the last lesson of chapter 447: I can use personal services: the hairdresser, blow-drying, the optician, the watchmaker and the jeweller.
\end{chapteropening}

\section[der Friseur]{der Friseur — the hairdresser, the barber}

\label{lesson:GE-C447-friseur}



\begin{quote}
Before the new one: say the German for to patch, then the German for to dry.
\end{quote}

\subsection*{You'll want to know: der Friseur}
\textbf{der Friseur} — \textquotedblleft{}the hairdresser, the barber\textquotedblright{}.

\begin{quote}
Einmal im Monat geht er zum Friseur.
\end{quote}

\begin{quote}
\emph{He goes to the hairdresser once a month.}
\end{quote}

Plural: \emph{die Friseure}. A woman: \emph{die Friseurin}. Also spelled \emph{Frisör}.

\begin{cousinweb}
\emph{Friseur} was borrowed from French \emph{friseur} 'hairdresser', from \emph{friser} 'to curl, to frizz' (EtymDB); French today mostly says \emph{coiffeur}.
\end{cousinweb}

\begin{grammarlens}[title={What you've built}]
One more word for telling, explaining and describing at B1.
\end{grammarlens}

\subsection*{Your turn}
\begin{itemize}
\raggedright
  \item \emph{Say it:} \emph{der Friseur}
  \item \emph{Say it:} \emph{der Friseur}, once more
  \item \emph{Say it:} a sentence of your own with \emph{der Friseur}
  \item \emph{Recall:} say the German for to patch, then the German for to dry, then say \emph{der Friseur} again
\end{itemize}

\begin{tcolorbox}[breakable,colback=teal!4,colframe=teal!35!black,title={Before you move on}]
What does der Friseur mean? (\textquotedblleft{}The hairdresser, the barber\textquotedblright{}.) Say it once more.
\end{tcolorbox}
\section[föhnen]{föhnen — to blow-dry}

\label{lesson:GE-C447-foehnen}



\begin{quote}
Before the new one: say the German for to dry, then the German for the hairdresser.
\end{quote}

\subsection*{You'll want to know: föhnen}
\textbf{föhnen} — \textquotedblleft{}to blow-dry\textquotedblright{}.

\begin{quote}
Nach dem Waschen föhnt sie sich die Haare.
\end{quote}

\begin{quote}
\emph{After washing her hair she blow-dries it.}
\end{quote}

\emph{Der Föhn} — the hair dryer, and also a warm wind from the mountains.

\begin{cousinweb}
\emph{föhnen} is built on \emph{Föhn}, a warm dry wind in the mountains, which goes back through Old High German \emph{phōnno} to Latin \emph{favōnius} 'mild west wind' (EtymDB; Lewis \& Short). The hair dryer was named after the wind.
\end{cousinweb}

\begin{grammarlens}[title={What you've built}]
One more word for telling, explaining and describing at B1.
\end{grammarlens}

\subsection*{Your turn}
\begin{itemize}
\raggedright
  \item \emph{Say it:} \emph{föhnen}
  \item \emph{Say it:} \emph{föhnen}, once more
  \item \emph{Say it:} a sentence of your own with \emph{föhnen}
  \item \emph{Recall:} say the German for to dry, then the German for the hairdresser, then say \emph{föhnen} again
\end{itemize}

\begin{tcolorbox}[breakable,colback=teal!4,colframe=teal!35!black,title={Before you move on}]
What does föhnen mean? (\textquotedblleft{}To blow-dry\textquotedblright{}.) Say it once more.
\end{tcolorbox}
\section[der Optiker]{der Optiker — the optician}

\label{lesson:GE-C447-optiker}



\begin{quote}
Before the new one: say the German for the hairdresser, then the German for to blow-dry.
\end{quote}

\subsection*{You'll want to know: der Optiker}
\textbf{der Optiker} — \textquotedblleft{}the optician\textquotedblright{}.

\begin{quote}
Beim Optiker sucht er eine neue Brille aus.
\end{quote}

\begin{quote}
\emph{At the optician's he chooses new glasses.}
\end{quote}

Same form in the plural. A woman: \emph{die Optikerin}.

\begin{cousinweb}
\emph{Optiker} is built on \emph{Optik} 'optics', a learned word from Greek \emph{optikós} 'of sight'. EtymDB does not list the German word.
\end{cousinweb}

\begin{grammarlens}[title={What you've built}]
One more word for telling, explaining and describing at B1.
\end{grammarlens}

\subsection*{Your turn}
\begin{itemize}
\raggedright
  \item \emph{Say it:} \emph{der Optiker}
  \item \emph{Say it:} \emph{der Optiker}, once more
  \item \emph{Say it:} a sentence of your own with \emph{der Optiker}
  \item \emph{Recall:} say the German for the hairdresser, then the German for to blow-dry, then say \emph{der Optiker} again
\end{itemize}

\begin{tcolorbox}[breakable,colback=teal!4,colframe=teal!35!black,title={Before you move on}]
What does der Optiker mean? (\textquotedblleft{}The optician\textquotedblright{}.) Say it once more.
\end{tcolorbox}
\section[der Uhrmacher]{der Uhrmacher — the watchmaker, the clockmaker}

\label{lesson:GE-C447-uhrmacher}



\begin{quote}
Before the new one: say the German for to blow-dry, then the German for the optician.
\end{quote}

\subsection*{You'll want to know: der Uhrmacher}
\textbf{der Uhrmacher} — \textquotedblleft{}the watchmaker, the clockmaker\textquotedblright{}.

\begin{quote}
Der Uhrmacher wechselt die Batterie in meiner Uhr.
\end{quote}

\begin{quote}
\emph{The watchmaker changes the battery in my watch.}
\end{quote}

Same form in the plural. A woman: \emph{die Uhrmacherin}.

\begin{cousinweb}
\emph{Uhrmacher} is built from \emph{Uhr} 'clock', which goes back to Latin \emph{hōra} 'hour' (EtymDB), and \emph{Macher} 'maker', from \emph{machen} 'to make' (EtymDB).
\end{cousinweb}

\begin{grammarlens}[title={What you've built}]
One more word for telling, explaining and describing at B1.
\end{grammarlens}

\subsection*{Your turn}
\begin{itemize}
\raggedright
  \item \emph{Say it:} \emph{der Uhrmacher}
  \item \emph{Say it:} \emph{der Uhrmacher}, once more
  \item \emph{Say it:} a sentence of your own with \emph{der Uhrmacher}
  \item \emph{Recall:} say the German for to blow-dry, then the German for the optician, then say \emph{der Uhrmacher} again
\end{itemize}

\begin{tcolorbox}[breakable,colback=teal!4,colframe=teal!35!black,title={Before you move on}]
What does der Uhrmacher mean? (\textquotedblleft{}The watchmaker, the clockmaker\textquotedblright{}.) Say it once more.
\end{tcolorbox}
\section[der Juwelier]{der Juwelier — the jeweller (US: jeweler)}

\label{lesson:GE-C447-juwelier}



\begin{quote}
Before the new one: say the German for the optician, then the German for the watchmaker.
\end{quote}

\subsection*{You'll want to know: der Juwelier}
\textbf{der Juwelier} — \textquotedblleft{}the jeweller (US: jeweler)\textquotedblright{}.

\begin{quote}
Beim Juwelier kaufen sie die Ringe für die Hochzeit.
\end{quote}

\begin{quote}
\emph{They buy the wedding rings at the jeweller's.}
\end{quote}

Plural: \emph{die Juweliere}. A woman: \emph{die Juwelierin}.

\begin{cousinweb}
\emph{Juwelier} is built on \emph{Juwel} 'jewel', which came from Old French \emph{joiel} (EtymDB), usually traced to Vulgar Latin \emph{jocāle}, built on Latin \emph{jocus} 'joke, play'.
\end{cousinweb}

\begin{grammarlens}[title={What you've built}]
One more word for telling, explaining and describing at B1.
\end{grammarlens}

\subsection*{Your turn}
\begin{itemize}
\raggedright
  \item \emph{Say it:} \emph{der Juwelier}
  \item \emph{Say it:} \emph{der Juwelier}, once more
  \item \emph{Say it:} a sentence of your own with \emph{der Juwelier}
  \item \emph{Recall:} say the German for the optician, then the German for the watchmaker, then say \emph{der Juwelier} again
\end{itemize}

\begin{tcolorbox}[breakable,colback=teal!4,colframe=teal!35!black,title={Before you move on}]
What does der Juwelier mean? (\textquotedblleft{}The jeweller (US: jeweler)\textquotedblright{}.) Say it once more.
\end{tcolorbox}
